% GENERATED FILE. Edit canonical lessons, then run npm run generate:books.
% canonical-source-hash: fnv1a64:ba6dd415bc646411
% canonical-lessons: FA-C197-shart, FA-C197-qanun, FA-C197-ba-in-hal, FA-C197-agarche, FA-C197-ya

\chapter{Condition, Law, Nevertheless, Although, Or}
\label{ch:fa-condition-law-nevertheless-although-or}
\hlchaptermodality{197}

\begin{chapteropening}
\textbf{By the end of this chapter:} \emph{I can say a condition, a law, nevertheless, although and or.}

Complete the last lesson of chapter 197: I can say a condition, a law, nevertheless, although and or.
\end{chapteropening}

\section[\texorpdfstring{\fa{شرط} (shart)}{شرط (shart)}]{\fa{شرط} (shart) — a condition}

\label{lesson:FA-C197-shart}



\begin{quote}
Before the new one: say the Persian for an edge, then the Persian for the ground floor.
\end{quote}

\subsection*{You'll want to know: \fa{شرط}}
\textbf{\fa{شرط}} — \emph{shart} — \textquotedblleft{}a condition\textquotedblright{}.

\begin{cousinweb}
\emph{shart} is borrowed from Arabic \emph{šarṭ}, \textquotedblleft{}condition, stipulation\textquotedblright{}. Its Arabic root is \emph{š-r-ṭ} (Lane).
\end{cousinweb}

\begin{grammarlens}[title={What you've built}]
One more word for this chapter.
\end{grammarlens}

\subsection*{Your turn}
\begin{itemize}
\raggedright
  \item \emph{Say it:} \emph{shart}
  \item \emph{Say it:} \emph{shart}, once more
  \item \emph{Say it:} say \emph{shart}
  \item \emph{Recall:} say the Persian for an edge, then the Persian for the ground floor, then say \emph{shart} again
\end{itemize}

\begin{tcolorbox}[breakable,colback=teal!4,colframe=teal!35!black,title={Before you move on}]
What does \fa{شرط} mean? (\textquotedblleft{}A condition\textquotedblright{}.) Say it once more.
\end{tcolorbox}
\section[\texorpdfstring{\fa{قانون} (qânun)}{قانون (qânun)}]{\fa{قانون} (qânun) — a law, a rule}

\label{lesson:FA-C197-qanun}



\begin{quote}
Before the new one: say the Persian for the ground floor, then the Persian for a condition.
\end{quote}

\subsection*{You'll want to know: \fa{قانون}}
\textbf{\fa{قانون}} — \emph{qânun} — \textquotedblleft{}a law, a rule\textquotedblright{}.

\begin{cousinweb}
\emph{qânun} is borrowed from Arabic \emph{qānūn}, \textquotedblleft{}rule, law\textquotedblright{}, which goes back to Greek \emph{kanōn}, \textquotedblleft{}rule, measuring rod\textquotedblright{} (EtymDB), the word behind English \emph{canon}.
\end{cousinweb}

\begin{grammarlens}[title={What you've built}]
One more word for this chapter.
\end{grammarlens}

\subsection*{Your turn}
\begin{itemize}
\raggedright
  \item \emph{Say it:} \emph{qânun}
  \item \emph{Say it:} \emph{qânun}, once more
  \item \emph{Say it:} say \emph{qânun}
  \item \emph{Recall:} say the Persian for the ground floor, then the Persian for a condition, then say \emph{qânun} again
\end{itemize}

\begin{tcolorbox}[breakable,colback=teal!4,colframe=teal!35!black,title={Before you move on}]
What does \fa{قانون} mean? (\textquotedblleft{}A law, a rule\textquotedblright{}.) Say it once more.
\end{tcolorbox}
\section[\texorpdfstring{\fa{با} \fa{این} \fa{حال} (bâ in hâl)}{با این حال (bâ in hâl)}]{\fa{با} \fa{این} \fa{حال} (bâ in hâl) — nevertheless, even so}

\label{lesson:FA-C197-ba-in-hal}



\begin{quote}
Before the new one: say the Persian for a condition, then the Persian for a law.
\end{quote}

\subsection*{You'll want to know: \fa{با} \fa{این} \fa{حال}}
\textbf{\fa{با} \fa{این} \fa{حال}} — \emph{bâ in hâl} — \textquotedblleft{}nevertheless, even so\textquotedblright{}.

\begin{cousinweb}
\emph{bâ in hâl} is \textquotedblleft{}with this state\textquotedblright{}: \emph{bâ}, \textquotedblleft{}with\textquotedblright{}, and \emph{in}, \textquotedblleft{}this\textquotedblright{}, are native; \emph{hâl}, \textquotedblleft{}state\textquotedblright{}, is borrowed from Arabic \emph{ḥāl} (root \emph{ḥ-w-l}, Lane).
\end{cousinweb}

\begin{grammarlens}[title={What you've built}]
One more word for this chapter.
\end{grammarlens}

\subsection*{Your turn}
\begin{itemize}
\raggedright
  \item \emph{Say it:} \emph{bâ in hâl}
  \item \emph{Say it:} \emph{bâ in hâl}, once more
  \item \emph{Say it:} say \emph{bâ in hâl}
  \item \emph{Recall:} say the Persian for a condition, then the Persian for a law, then say \emph{bâ in hâl} again
\end{itemize}

\begin{tcolorbox}[breakable,colback=teal!4,colframe=teal!35!black,title={Before you move on}]
What does \fa{با} \fa{این} \fa{حال} mean? (\textquotedblleft{}Nevertheless, even so\textquotedblright{}.) Say it once more.
\end{tcolorbox}
\section[\texorpdfstring{\fa{اگرچه} (agarche)}{اگرچه (agarche)}]{\fa{اگرچه} (agarche) — although}

\label{lesson:FA-C197-agarche}



\begin{quote}
Before the new one: say the Persian for a law, then the Persian for nevertheless.
\end{quote}

\subsection*{You'll want to know: \fa{اگرچه}}
\textbf{\fa{اگرچه}} — \emph{agarche} — \textquotedblleft{}although\textquotedblright{}.

\begin{cousinweb}
\emph{agarche} is \emph{agar}, \textquotedblleft{}if\textquotedblright{}, with \emph{che}, \textquotedblleft{}what\textquotedblright{}: \textquotedblleft{}even if\textquotedblright{}. Both are native; \emph{agar} continues Middle Persian \emph{agar} (EtymDB).
\end{cousinweb}

\begin{grammarlens}[title={What you've built}]
One more word for this chapter.
\end{grammarlens}

\subsection*{Your turn}
\begin{itemize}
\raggedright
  \item \emph{Say it:} \emph{agarche}
  \item \emph{Say it:} \emph{agarche}, once more
  \item \emph{Say it:} say \emph{agarche}
  \item \emph{Recall:} say the Persian for a law, then the Persian for nevertheless, then say \emph{agarche} again
\end{itemize}

\begin{tcolorbox}[breakable,colback=teal!4,colframe=teal!35!black,title={Before you move on}]
What does \fa{اگرچه} mean? (\textquotedblleft{}Although\textquotedblright{}.) Say it once more.
\end{tcolorbox}
\section[\texorpdfstring{\fa{یا} (yâ)}{یا (yâ)}]{\fa{یا} (yâ) — or}

\label{lesson:FA-C197-ya}



\begin{quote}
Before the new one: say the Persian for nevertheless, then the Persian for although.
\end{quote}

\subsection*{You'll want to know: \fa{یا}}
\textbf{\fa{یا}} — \emph{yâ} — \textquotedblleft{}or\textquotedblright{}.

\begin{cousinweb}
\emph{yâ} is a native Persian word, inherited rather than borrowed; it is usually traced to Middle Persian \emph{ayāb}, \textquotedblleft{}or\textquotedblright{}.
\end{cousinweb}

\begin{grammarlens}[title={What you've built}]
One more word for this chapter.
\end{grammarlens}

\subsection*{Your turn}
\begin{itemize}
\raggedright
  \item \emph{Say it:} \emph{yâ}
  \item \emph{Say it:} \emph{yâ}, once more
  \item \emph{Say it:} say \emph{yâ}
  \item \emph{Recall:} say the Persian for nevertheless, then the Persian for although, then say \emph{yâ} again
\end{itemize}

\begin{tcolorbox}[breakable,colback=teal!4,colframe=teal!35!black,title={Before you move on}]
What does \fa{یا} mean? (\textquotedblleft{}Or\textquotedblright{}.) Say it once more.
\end{tcolorbox}
